\chapter{Guía de arranque del sistema}
\label{cap:AnexoK}

Este anexo describe los pasos necesarios para poner en marcha el sistema en un entorno de desarrollo local. Toda la infraestructura se gestiona mediante Docker Compose y un \texttt{Makefile} que centraliza los comandos habituales.

\section{Requisitos previos}

\begin{itemize}
 \item \textbf{Docker} 24 o superior con el plugin \texttt{docker compose} (v2).
 \item \textbf{GNU Make}.
 \item \textbf{Git}.
 \item \textit{Opcional, modo GPU:} controlador NVIDIA y \texttt{nvidia-container-toolkit} instalados en el sistema anfitrión.
\end{itemize}

\section{Configuración del entorno}

Antes de arrancar por primera vez es necesario crear el fichero de variables de entorno a partir de la plantilla incluida en el repositorio:

\begin{lstlisting}[language=bash]
cp .env.example .env
\end{lstlisting}

Los valores por defecto son válidos para un entorno local de desarrollo. El catálogo completo de variables, con sus valores por defecto y el servicio que las consume, se recoge en el Anexo~\ref{cap:AnexoB} (Tabla~\ref{tab:envvars}).

\section{Primera puesta en marcha}

La primera vez es necesario ejecutar dos comandos en orden. El primero prepara el entorno: detiene los servicios previos y elimina sus volúmenes (incluida la base de datos), construye las imágenes Docker, descarga el modelo desde Hugging Face (\texttt{make model-download}, que usa \texttt{HUGGING\_FACE\_HUB\_TOKEN} del \texttt{.env}), crea y migra las bases de datos y carga los datos iniciales y el catálogo CIE-10:

\begin{lstlisting}[language=bash]
make setup
\end{lstlisting}

Este proceso puede tardar varios minutos la primera vez, principalmente por la descarga de imágenes base y la compilación de dependencias Elixir. \textbf{No levanta los servicios}; una vez completado, hay que arrancarlos explícitamente:

\begin{lstlisting}[language=bash]
make up
\end{lstlisting}

\section{Arranque y parada habituales}

A partir de la segunda vez, el ciclo de trabajo diario se reduce a dos comandos (\texttt{make up} detecta la GPU con \texttt{nvidia-smi} y, si no hay, arranca el motor de IA en CPU):

\begin{lstlisting}[language=bash]
make up # Levanta todos los servicios en segundo plano
make down # Detiene todos los servicios
\end{lstlisting}

Tras ejecutar \texttt{make up} los servicios quedan accesibles en las siguientes URLs (tabla~\ref{tab:urls}).

\begin{table}[h]
\centering
\caption{URLs de acceso a los servicios en entorno local.}
\label{tab:urls}
\begin{tabular}{ll}
\hline
\textbf{Servicio} & \textbf{URL} \\
\hline
Frontend (Next.js) & \url{http://localhost:3000} \\
Backend API (Phoenix) & \url{http://localhost:4000} \\
Motor de IA (FastAPI) & \url{http://localhost:8000} \\
Backoffice admin (Phoenix) & \url{http://localhost:4000/admin} \\
Mailpit (bandeja SMTP, solo desarrollo) & \url{http://localhost:8025} \\
Prometheus & \url{http://localhost:9090} \\
Grafana & \url{http://localhost:3030} \\
PostgreSQL & \texttt{localhost:5432} \\
\hline
\end{tabular}
\\[4pt]
{\footnotesize Fuente: objetivo \texttt{up} del \texttt{Makefile}.}
\end{table}

\section{Modos de ejecución}

El sistema admite tres modos de ejecución en función de la disponibilidad de hardware y del estado del modelo entrenado (tabla~\ref{tab:modos}).

\begin{table}[h]
\centering
\caption{Modos de ejecución disponibles.}
\label{tab:modos}
\begin{tabular}{llP{6.2cm}}
\hline
\textbf{Modo} & \textbf{Comando} & \textbf{Requisito} \\
\hline
GPU & \texttt{make up} (GPU detectada) & GPU NVIDIA + modelo entrenado (preferente, entrenamiento/desarrollo/producción) \\
CPU & \texttt{make cpu-up} & Modelo entrenado (soporte, cuando no hay GPU disponible) \\
Mock & \texttt{make mock-up} & Ninguno (respuestas simuladas) \\
\hline
\end{tabular}
\\[4pt]
{\footnotesize Fuente: objetivos \texttt{up} (con la GPU detectada), \texttt{cpu-up} y \texttt{mock-up} del \texttt{Makefile}.}
\end{table}

Cada modo se describe a continuación:

\begin{sloppypar}
\begin{itemize}
 \item \textbf{GPU} (\texttt{make up}, con GPU detectada): requiere una GPU NVIDIA con \texttt{nvidia-container-toolkit} y un modelo previamente entrenado montado en \texttt{ai\_engine/model/}. Es el modo preferente siempre que hay GPU disponible, tanto en desarrollo como en producción (\texttt{make deploy GPU=1}, Anexo~\ref{cap:AnexoI}) y durante el entrenamiento, porque ofrece la máxima velocidad de inferencia.
 \item \textbf{CPU} (\texttt{make cpu-up}): carga el modelo real pero realiza la inferencia en CPU. Es el modo de soporte para cuando no hay GPU disponible, sea en desarrollo o en un despliegue sin acelerador (Anexo~\ref{cap:AnexoI}): proponer los códigos cuesta 0{,}331\,s por informe, aunque justificarlos sube a 16{,}2\,s, y de ahí que la atribución se sirva en una petición aparte.
 \item \textbf{Mock} (\texttt{make mock-up}): el motor de IA devuelve respuestas fijas sin cargar ningún modelo. Recomendado para desarrollo del frontend y pruebas funcionales del backend cuando no se dispone de modelo entrenado.
\end{itemize}
\end{sloppypar}

\section{Usuarios de prueba}

El comando \texttt{make setup} (a través de \texttt{mix run priv/repo/seeds.exs}) crea automáticamente el siguiente usuario administrador con la cuenta ya confirmada (tabla~\ref{tab:seeds}).

\begin{table}[h]
\centering
\caption{Usuario de prueba creado por el seed.}
\label{tab:seeds}
\begin{tabular}{llll}
\hline
\textbf{Email} & \textbf{Contraseña} & \textbf{Username} & \textbf{Rol} \\
\hline
\texttt{myadminuser@clicoder.app} & \texttt{password123!} & \texttt{admin} & Administrador \\
\hline
\end{tabular}
\\[4pt]
{\footnotesize Fuente: \texttt{backend/\allowbreak{}priv/\allowbreak{}repo/\allowbreak{}seeds.exs}.}
\end{table}

El correo y la contraseña del administrador son configurables mediante las variables de entorno \texttt{SEED\_ADMIN\_EMAIL} y \texttt{SEED\_ADMIN\_PASSWORD} (los valores de la tabla son los predeterminados). El usuario administrador tiene acceso al backoffice en \url{http://localhost:4000/admin}.

\section{Consulta de logs}

Para inspeccionar la salida de cualquier servicio en tiempo real:

\begin{lstlisting}[language=bash]
make logs # Menú interactivo para seleccionar el servicio
\end{lstlisting}

\section{Reset de base de datos}

Si es necesario partir de un estado limpio sin reconstruir imágenes:

\begin{lstlisting}[language=bash]
make db-reset # Drop + migrate + seed
\end{lstlisting}

\section{Motor de IA: diccionario y entrenamiento}

El motor de IA opera en dos modos complementarios: un clasificador neuronal basado en RigoBERTa-Clinical y un clasificador por diccionario (\emph{baseline}) basado en coincidencia de patrones. Ambos requieren pasos de preparación independientes.

\subsection{Regenerar el diccionario de términos clínicos}

El clasificador por diccionario utiliza un fichero \texttt{model/baseline\_dict.json} que mapea términos clínicos lematizados a bloques CIE-10. Este fichero se construye con tres fuentes, las que pasa el \texttt{Makefile} en \texttt{--sources clinical corpus combined}: \texttt{clinical} (términos clínicos reales extraídos de las anotaciones de la tarea~X de CodiEsp, ampliados con sinónimos léxicos), \texttt{corpus} (\emph{n}-gramas discriminativos minados de las notas de entrenamiento) y \texttt{combined}, la unión de las dos anteriores, que es el diccionario que se guarda. Las descripciones oficiales del catálogo (\texttt{diagnoses}, \texttt{procedures}, \texttt{chemicals}) existen como fuentes opcionales de \texttt{baseline\_dict.py}, pero no entran en el diccionario servido porque, medidas, empeoran el resultado. Adicionalmente, el sistema expande abreviaturas clínicas frecuentes en español (\texttt{HTA}, \texttt{DM2}, \texttt{EPOC}, \texttt{IAM}\dots) a su forma completa antes de lematizar, siguiendo la estrategia del equipo IAM ganador de CodiEsp~2020.

Es necesario regenerarlo cuando se modifica \texttt{baseline\_dict.py}, se actualizan los ficheros de datos de CodiEsp o se cambia cualquier parámetro de construcción del diccionario. El proceso lee los datos desde \texttt{/data/}, lematiza todos los términos con spaCy (el resultado se almacena en caché en \texttt{model/baseline\_cache/} para acelerar ejecuciones sucesivas) y escribe el diccionario resultante:

\begin{lstlisting}[language=bash]
make ai-baseline-dict 2>&1 | tee /tmp/dict.log
\end{lstlisting}

La primera ejecución puede tardar entre 15 y 30~minutos porque lematiza las decenas de miles de entradas del catálogo y las notas del corpus. Las ejecuciones posteriores son mucho más rápidas gracias a la caché. Al finalizar se imprime un resumen con el número de bloques CIE-10 cubiertos y los patrones generados por cada fuente.

Para comparar el rendimiento con y sin expansión de abreviaturas (útil para medir el impacto aislado de esa mejora), se puede desactivar puntualmente:

\begin{lstlisting}[language=bash]
make ai-baseline-dict NO_ABBREVS=1 2>&1 | tee /tmp/dict_noabbrev.log
\end{lstlisting}

\subsection{Entrenar el clasificador neuronal}

El clasificador neuronal se entrena con \texttt{train.py} sobre el corpus CodiEsp. El proceso descarga el modelo base desde Hugging Face (requiere \texttt{HF\_TOKEN} en \texttt{.env}), ajusta mediante \emph{fine-tuning} el \emph{encoder} RigoBERTa-Clinical con una cabeza de clasificación \emph{multi-label} y guarda el modelo resultante en \texttt{ai\_engine/model/}. El objetivo \texttt{make ai-train-gpu} lanza \texttt{train.py} con GPU explícita (\texttt{--device cuda}) a través del servicio \texttt{ai\_engine} y el fichero \texttt{docker-compose.gpu.yml}. Exige que la imagen del servicio se haya construido con soporte CUDA (\texttt{make build-ai GPU=1}): con una imagen construida solo para CPU falla con «Torch not compiled with CUDA enabled». En el equipo de desarrollo el entrenamiento se lanzó con \texttt{train.py} directamente sobre otra imagen que ya incluye PyTorch con CUDA (\texttt{cie-10-training}). Ejecutado sin variables usa los valores por defecto del \texttt{Makefile} (\texttt{MAX\_LENGTH=1024}, \texttt{LR=5e-6}, \texttt{DROPOUT=0.1}, \texttt{WEIGHT\_DECAY=0.01}, \texttt{FREEZE\_LAYERS=0}, \texttt{EPOCHS=20}), que \textbf{no} reproducen el modelo servido. Para reproducirlo hay que pasar todos los hiperparámetros de la ejecución \texttt{20260921T091919Z} registrada en \texttt{training\_runs.csv}:

\begin{lstlisting}[language=bash, breaklines=true, breakatwhitespace=false]
make ai-train-gpu MAX_LENGTH=512 BATCH_SIZE=8 GRAD_ACCUM=4 \
  THRESHOLD=0.3 LR=1e-4 WARMUP_RATIO=0.1 WEIGHT_DECAY=0.1 \
  DROPOUT=0.3 FREEZE_LAYERS=20 EPOCHS=150 PATIENCE=20 \
  FULL_CODES=1 LABEL_SMOOTHING=0.1 LR_SCHEDULE=plateau \
  UNFREEZE_EVERY=30 UNFREEZE_LAYERS=4 UNFREEZE_LR_RATIO=0.1 \
  PRETRAIN_EPOCHS=10 LAMBDA_HIER=0.5 SELECT_METRIC=map_macro \
  RANK_LOSS_WEIGHT=0.1 2>&1 | tee /tmp/train.log
\end{lstlisting}

El preentrenamiento con los CSV de la tarea X de CodiEsp es opcional y se activa con \texttt{PRETRAIN\_TASKX=<ruta>}; el comando anterior no lo usa. El resto de parámetros (\texttt{POS\_WEIGHT\_CAP=10}, \texttt{CHUNK\_OVERLAP=64}, \texttt{ASL\_CLIP=0.05}) coinciden con los valores por defecto del \texttt{Makefile} y de \texttt{train.py} (\texttt{ASL\_CLIP} solo está definido en este último); la ejecución original no fijó semilla, por lo que una repetición no reproduce las métricas al decimal. Este mismo comando, con las variables adicionales que se indiquen (por ejemplo \texttt{TRAIN\_FILE}), es el que se toma como base en el Anexo~\ref{cap:AnexoL}.

Cada ejecución guarda un \emph{checkpoint} con el mejor valor de validación de la métrica de selección (\texttt{SELECT\_METRIC}; el comando anterior usa \texttt{map\_macro}), los umbrales por clase optimizados y un registro de métricas por época en \texttt{model/training\_runs.csv}. La ejecución del modelo servido duró 1.409~s (unos 23,5~minutos) en la RTX 4080, y la más larga de las 58 ejecuciones registradas sobre 1.767 códigos, 1,7~horas (columna \texttt{total\_seconds} de \texttt{training\_runs.csv}).

\section{Monitorización con Prometheus y Grafana}

El sistema incluye una pila de observabilidad compuesta por Prometheus y Grafana, que en desarrollo se levanta con \texttt{make start-monitoring-dev} (\texttt{make up} no la arranca, aunque muestra sus URLs).

\subsection{URLs y credenciales}

Las URLs y credenciales de acceso están en la Tabla~\ref{tab:monitoring-urls}.

\begin{table}[h]
\centering
\caption{Acceso a las herramientas de monitorización en entorno local.}
\label{tab:monitoring-urls}
\begin{tabular}{lll}
\hline
\textbf{Herramienta} & \textbf{URL} & \textbf{Credenciales} \\
\hline
Prometheus & \url{http://localhost:9090} & Sin autenticación \\
Grafana & \url{http://localhost:3030} & \texttt{admin} / \texttt{GRAFANA\_PASSWORD} \\
\hline
\end{tabular}
\\[4pt]
{\footnotesize Fuente: \texttt{docker-compose.\allowbreak{}monitoring.\allowbreak{}dev.yml}.}
\end{table}

La contraseña de Grafana se configura mediante la variable \texttt{GRAFANA\_PASSWORD} del fichero \texttt{.env} (valor por defecto: \texttt{changeme}). La fuente de datos de Prometheus queda provisionada automáticamente al arrancar: no es necesario configurarla manualmente en la interfaz de Grafana.

\subsection{Métricas expuestas}

Prometheus recoge métricas de ocho tareas de recogida (\emph{jobs}) con un intervalo de muestreo de 15 segundos (tabla~\ref{tab:prometheus-jobs}).

\begin{table}[h]
\centering
\caption{Tareas de recogida (\emph{jobs}) configuradas en Prometheus.}
\label{tab:prometheus-jobs}
\begin{tabular}{llP{6.2cm}}
\hline
\textbf{Job} & \textbf{Endpoint} & \textbf{Métricas principales} \\
\hline
\texttt{backend} & \texttt{:9568/metrics} & Latencia HTTP (Phoenix), tiempos de consulta Ecto, memoria BEAM, colas de planificadores \\
\texttt{ai\_engine} & \texttt{:8000/metrics} & Latencia de inferencia, peticiones totales, estado del modelo \\
\texttt{prometheus} & \texttt{:9090/metrics} & Métricas internas de Prometheus \\
\texttt{cadvisor} & \texttt{:8080/metrics} & CPU, memoria, red y disco por contenedor \\
\texttt{node\_exporter} & \texttt{:9100/metrics} & Métricas del host: CPU, RAM, disco y carga \\
\texttt{traefik} & \texttt{:8082/metrics} & Peticiones, códigos de estado y latencia del proxy inverso \\
\texttt{cloudflared} & \texttt{:20241/metrics} & Estado y tráfico del túnel de Cloudflare \\
\texttt{blackbox} & \texttt{:9115/probe} & Sondas de disponibilidad HTTP de los servicios públicos \\
\hline
\end{tabular}
\\[4pt]
{\footnotesize Fuente: \texttt{docker/\allowbreak{}prometheus/\allowbreak{}prometheus.yml}.}
\end{table}

El backend Elixir utiliza \texttt{telemetry\_metrics\_prometheus} para exponer las métricas en el puerto \texttt{9568}. Las métricas de VM BEAM configuradas son \texttt{vm.memory.total} (memoria total del runtime) y \texttt{vm.total\_run\_queue\_lengths} (longitud de las colas de planificadores de CPU y de E/S), útiles para detectar saturación en el procesamiento concurrente de WebSockets. El motor de IA expone las suyas en el endpoint \texttt{/metrics} del puerto \texttt{8000} mediante \texttt{prometheus\_fastapi\_instrumentator} (métricas HTTP automáticas) más dos métricas propias: \texttt{cie10\_inference\_duration\_seconds} (histograma de latencia de inferencia) y los indicadores (\emph{gauges}) \texttt{cie10\_model\_loaded} y \texttt{cie10\_dict\_loaded} (estado de carga del modelo BERT y del clasificador de diccionario).

\subsection{Retención de datos}

Prometheus almacena los datos de series temporales durante \textbf{15 días} (configurado con \texttt{--storage.tsdb.retention.time=15d}). Los datos persisten entre reinicios del contenedor gracias al volumen Docker \texttt{prometheus\_data}. Los paneles de Grafana persisten en el volumen \texttt{grafana\_data}.

\subsection{Exploración de métricas}

Para explorar las métricas disponibles sin necesidad de construir un panel en Grafana, se puede usar directamente la interfaz web de Prometheus en \url{http://localhost:9090/graph}. Ejemplos de consultas PromQL útiles:

\begin{lstlisting}[language=bash]
# Latencia de inferencia del motor de IA (percentil 95)
histogram_quantile(0.95, rate(cie10_inference_duration_seconds_bucket[5m]))

# Estado de carga del modelo BERT (1 = cargado, 0 = no cargado)
cie10_model_loaded

# Memoria total de la VM BEAM en kilobytes
vm_memory_total_kilobytes

# Longitud de la cola de planificadores de CPU del runtime Elixir
vm_total_run_queue_lengths_cpu
\end{lstlisting}
